\documentclass[12pt,a4paper]{article}
\usepackage[utf8]{inputenc}
\usepackage[T2A]{fontenc}
\usepackage[russian]{babel}
\usepackage{amsmath,amssymb,amsthm}
\usepackage{booktabs,array}
\usepackage{geometry}
\usepackage{microtype}
\geometry{margin=2.3cm}

\newtheorem{proposition}{Предложение}
\newtheorem{protocol}{Протокол}
\newtheorem{nogocriterion}{Критерий провала}

\title{SU(5)-селектор рангов для $\alpha_s$ и угла Вайнберга}
\author{}
\date{4 августа 2026 г.}

\begin{document}
\maketitle

\section{Каноническое разложение пространства размерности 24}

Пусть
\begin{equation}
 P=\operatorname{diag}(1,1,1,-1,-1)
\end{equation}
задаёт $\mathbb Z_2$-инволюцию в $SU(5)$, а гиперзаряд нормирован как
\begin{equation}
 Y=\operatorname{diag}\left(-\frac13,-\frac13,-\frac13,
 \frac12,\frac12\right).
\end{equation}
Тогда комплексированное присоединённое представление раскладывается в
\begin{equation}
 \mathbf{24}
 =({\bf8},{\bf1})_0
 \oplus({\bf1},{\bf3})_0
 \oplus({\bf1},{\bf1})_0
 \oplus({\bf3},{\bf2})_{-5/6}
 \oplus(\overline{\bf3},{\bf2})_{+5/6}.
\end{equation}
Размерности блоков равны
\begin{equation}
 24=8+3+1+6+6.
\end{equation}
Первые три блока имеют чётность $P=+1$, два смешанных блока --- $P=-1$;
получается точное разбиение $12+12$.

\section{Сильная связь}

Для атомарного проектора $Q$ введём минимальный функционал
\begin{equation}
 \alpha_s(Q)
 =\frac{1}{\operatorname{rank}Q+pi^2
 \operatorname{rank}Q/24}.
\end{equation}
На рангах $8,3,1,6,6$ только смешанная орбита $X/\overline X$ попадает в
округлённое атласом значение $0.1181$:
\begin{equation}
 \alpha_s(X)
 =\frac{1}{6+\pi^2/4}
 =0.1180999917\ldots
\end{equation}
Здесь число $6$ и коэффициент $1/4$ не выбираются независимо: они являются
обычным и нормированным следами одного rank-$6$ проектора.

\section{Угол Вайнберга}

Обозначим ранги цветового, слабого, гиперзарядного и смешанного блоков через
$r_C,r_W,r_Y,r_X$. Атласная формула допускает rank-запись
\begin{equation}
 \sin^2\theta_W
 =\frac{r_C-r_W(r_X/24)\pi^{-1}}
 {(24-r_W)+(r_W+r_Y)\pi}.
\end{equation}
Каноническое назначение
\begin{equation}
 (r_C,r_W,r_Y,r_X)=(8,3,1,6)
\end{equation}
даёт
\begin{equation}
 \sin^2\theta_W
 =\frac{8-3/(4\pi)}{21+4\pi}
 =0.2312215305\ldots
\end{equation}

Проверены все $120$ назначения четырёх различных ролей пяти блокам. Физическое
назначение занимает первое место; единственная ничья получается при замене
$X\leftrightarrow\overline X$. После факторизации зарядового сопряжения
победитель единственен. Следующий неэквивалентный вариант даёт $0.2226941$ и
не попадает в интервал округления $0.2312\pm0.00005$.

\begin{proposition}[Rank-selector gate]
Все целые и дробные коэффициенты двух атласных формул восстанавливаются из
канонических рангов блоков присоединённого представления $SU(5)$. В полном
перестановочном контроле оба физических назначения уникальны с точностью до
зарядового сопряжения.
\end{proposition}

\section{Что найдено и чего ещё нет}

Результат существенно сильнее наблюдения ``дроби делят 24'': один rank-$6$
проектор одновременно объясняет $6$ и $1/4$ в $\alpha_s$, а формула Вайнберга
использует именно физические блоки $C,W,Y,X$ в их физических ролях.

Однако функциональные формы были прочитаны из атласа до rank-реконструкции.
Представления объясняют коэффициенты, но пока не объясняют, почему действие
должно собрать их именно как
\begin{equation}
 r+\pi^2r/24
 \quad\text{и}\quad
 \frac{r_C-r_W(r_X/24)\pi^{-1}}
 {(24-r_W)+(r_W+r_Y)\pi}.
\end{equation}

\begin{nogocriterion}[Граница повышения статуса]
Текущий результат является уникальной representation-theoretic
реконструкцией, но не blind-предсказанием. Для повышения статуса оба
функционала должны следовать из одного заранее заданного gauge-fixed action.
После этого то же правило обязано вычислить скрытую третью связь или порог без
изменения грамматики.
\end{nogocriterion}

Машинный аудит сохранён в
\texttt{s2t\_su5\_rank\_selector\_results.json} и воспроизводится программой
\texttt{s2t\_su5\_rank\_selector\_audit.py}.

\end{document}